% Work-MAT -- 2. rocnik, hodina 09. Metodicka priprava, nikoli zadani testu.
\documentclass[11pt,a4paper]{article}
\input{../../sablony/pripravy/priprava.sty}
\nastavPripravu{2. ročník}{Inverzní a mocninné funkce}{09}{Průběžný test T1 a rozbor}
\begin{document}
\poznamka{Toto je výhradně učitelská příprava. Neobsahuje schválené zadání testu ani bodování či klasifikační hranice. Před hodinou připrav konkrétní test z platné šablony a ověř jeho rozsah, varianty, řešení a tisk. Doporučené načasování: 3 min instrukce, 15--20 min test, 12 min anonymní rozbor typických postupů, 8 min krátké upevnění, zbytek organizační závěr.}

\blok{Před hodinou: kontrolní seznam}
\begin{itemize}[label=$\square$]
\item Ověřit, že test prověřuje pouze \textbf{probraný obsah hodin 03--08} a potřebné opakování z hodiny 02.
\item Připravit odpovídající varianty zadání a samostatná učitelská řešení; předem provést vlastní matematickou kontrolu.
\item Vycházet z aktuálních pravidel pro psaní testů, pomůcky, absence a náhradní termíny zveřejněných v Moodle.
\item Připravit čistý papír na výpočty a promyslet sběr testů; přesné pokyny sdělit žákům před začátkem.
\end{itemize}

\blok{Obsahový rámec T1}
Test má být krátký a přiměřený času. Nabízí se ověřit dvě až tři odlišné dovednosti, například:
\begin{itemize}
\item poznání prosté funkce a podmínek existence inverze;
\item získání inverzního předpisu včetně definičního oboru;
\item základní vlastnosti $x^n$ pro kladný a záporný celočíselný exponent, zejména definiční obor, paritu a obor hodnot.
\end{itemize}
\poznamka{Samotný test do této přípravy nevkládám. Právě tak zde neurčuji body, váhy ani hranice klasifikace. Hodina 10 o transformacích mocninných funkcí \textbf{ještě nenastala}, a proto nemá být jejím obsahem test T1 podmiňován.}

\blok{Po odevzdání: rozbor bez prozrazování variant}
Připomeň žákům rozdíl mezi prostotou a monotónností, mezi inverzí a převrácenou hodnotou funkce a mezi $x^2$ na $\mathbb R$ a na $\langle0;+\infty)$. Nezveřejňuj řešení testových úloh, pokud se tytéž varianty budou používat při dopisování; místo toho pracuj na jiných, analogických příkladech.

\begin{prikladReseniBox}[title={Příklad 1}]
Pro funkci $f(x)=2x-5$ na $\mathbb R$ určete předpis inverze a ověřte jej složením.
\tcblower
\[
\begin{aligned}
f(x)&=2x-5 &&\bigm|\ x\to y\\
x&=2y-5\\
2y&=x+5\\
y&=\frac{x+5}{2}\\
f^{-1}(x)&=\frac{x+5}{2}.
\end{aligned}
\]
Kontrola: $f^{-1}(f(x))=(2x-5+5)/2=x$ pro každé $x\in\mathbb R$.
\end{prikladReseniBox}

\begin{prikladReseniBox}[title={Příklad 2}]
Porovnej definiční obory, obory hodnot a paritu funkcí $g(x)=x^4$ a $h(x)=x^{-4}$.
\tcblower
$g$ má $D=\mathbb R$, $H=\langle0;+\infty)$ a je sudá; $h$ má $D=\mathbb R\setminus\{0\}$, $H=(0;+\infty)$ a je rovněž sudá. První funkce dosahuje nuly, druhá nikdy.
\end{prikladReseniBox}

\blok{Krátké nehodnocené upevnění}
\begin{zadaniReseniBox}[title={Samostatná práce 1}]
\begin{enumerate}[label=\alph*)]
\item Je funkce $f(x)=x^2$ na $\langle0;+\infty)$ prostá? Jaký má inverzní předpis?
\item Jaké jsou asymptoty grafu $g(x)=x^{-3}$?
\end{enumerate}
\tcblower
\begin{enumerate}[label=\alph*)]
\item Ano, protože na uvedeném oboru přísně roste.
\(\displaystyle\begin{aligned}f(x)&=x^2 &&\bigm|\ x\to y\\x&=y^2\\y&=\sqrt{x}\\f^{-1}(x)&=\sqrt{x}\end{aligned}\)

$D(f^{-1})=\langle0;+\infty)$.
\item $x=0$ a $y=0$.
\end{enumerate}
\end{zadaniReseniBox}

\poznamka{Pokud do konce hodiny není čas na upevnění, ulož je jako dobrovolné procvičení a věnuj se organizačnímu uzavření. Záznam o reálném průběhu a případných nepřítomných proveď podle školních pravidel.}

\begin{zadaniReseniBox}[title={Domácí úkol}]
Zapiš si dva pojmy nebo postupy z bloku inverzních a mocninných funkcí, ve kterých si stále nejsi jistý, a ke každému vymysli jednoduchý vlastní příklad. Nejprve zkus příklady vyřešit bez poznámek.
\tcblower
Jde o sebereflexi. Např. „prostota“: $x^2$ na $\mathbb R$ není prostá; „záporný exponent“: $x^{-2}=1/x^2$, proto $x\ne0$. Konkrétní příklady se budou mezi žáky lišit.
\end{zadaniReseniBox}
\end{document}
